\section{Methods}
\label{sec:method}

\subsection{Designs}
\label{sec:designs}

Every design answers \texttt{GET /result?roll=$n$} with the same JSON body and
gets the same budget of two CPUs (hardware threads) for answering students
(Fig.~\ref{fig:designs}). Before the announced release time every design
answers ``not released''.
\textbf{A} builds each page from PostgreSQL (a 4-table join).
\textbf{M} splits the budget into two independent copies of A (``mirror
links''); each student picks one at random and, on retry, switches with
probability 0.5.
\textbf{B0} puts a Redis cache in front of A (cache-aside); \textbf{B1} adds
request coalescing, \textbf{B2} probabilistic early
recomputation~\cite{vattani2015xfetch}, and \textbf{B3} fills the cache before
release.
\textbf{C1} pre-computes every result into Redis, and \textbf{C2} writes one
JSON file per student, served by Nginx; both keep an idle back-office
PostgreSQL (0.25 CPU, counted in cost).
\textbf{D} puts a waiting room (0.25 CPU) in front of A, C1 or C2 (we test
D-A and D-C1). It admits
students at 80\% of the backend's capacity, either first come, first served
(FIFO) or by a lottery among early arrivals, as some products
offer~\cite{cloudflare2026queue}. Its ticket is either \emph{naive} (held in
the page and lost when the student presses refresh) or a signed \emph{token}
cookie (kept on refresh). An admitted student may use the backend for
\SI{60}{s}; the waiting room limits neither how many admitted students are
active nor how many requests each sends.

\begin{figure*}[t]
% Read paths of the designs (Section III). Included by sections/designs.tex.
\tikzset{
  box/.style={draw, rounded corners=1.5pt, minimum height=7mm, minimum width=13mm,
              inner sep=2pt, align=center, font=\footnotesize, fill=white},
  stu/.style={box, fill=gray!12},
  app/.style={box, fill=blue!8},
  db/.style={draw, cylinder, shape border rotate=90, aspect=0.22, minimum height=8mm,
             minimum width=14mm, inner sep=1pt, align=center, font=\footnotesize, fill=orange!12},
  kv/.style={db, fill=red!10},
  files/.style={box, fill=green!10, double copy shadow={shadow xshift=1.2pt, shadow yshift=1.2pt}},
  wr/.style={box, fill=yellow!20},
  req/.style={-{Stealth[length=1.8mm]}, semithick},
  pre/.style={-{Stealth[length=1.8mm]}, semithick, dashed, gray!60!black},
  lab/.style={font=\scriptsize, inner sep=1pt, fill=white},
}
\newcommand{\panel}[2]{\begin{minipage}[b]{0.325\textwidth}\centering
  \begin{tikzpicture}[node distance=4mm and 6mm]#2\end{tikzpicture}\par\vspace{2pt}
  {\footnotesize #1}\end{minipage}}
\centering
\panel{(A) Database-backed: every request\\builds the page (+5\,ms CPU)}{%
  \node[stu] (s) {Students};
  \node[app, right=of s] (a) {App};
  \node[db, right=of a] (d) {Postgre-\\SQL};
  \draw[req] (s) -- (a); \draw[req] (a) -- (d);
}\hfill
\panel{(M) Mirrors: two copies of A,\\each with half the CPU}{%
  \node[stu] (s) {Students};
  \node[app, right=of s, yshift=5mm] (a1) {App 1};
  \node[app, right=of s, yshift=-5mm] (a2) {App 2};
  \node[db, right=of a1, minimum height=6.5mm] (d1) {DB 1};
  \node[db, right=of a2, minimum height=6.5mm] (d2) {DB 2};
  \draw[req] (s) -- (a1); \draw[req] (s) -- (a2);
  \draw[req] (a1) -- (d1); \draw[req] (a2) -- (d2);
}\hfill
\panel{(B) Cache-aside: Redis first,\\database only on a miss}{%
  \node[stu] (s) {Students};
  \node[app, right=of s] (a) {App};
  \node[kv, right=of a, yshift=6mm, minimum height=6.5mm] (r) {Redis};
  \node[db, right=of a, yshift=-6mm, minimum height=6.5mm] (d) {Postgre-\\SQL};
  \draw[req] (s) -- (a);
  \draw[req] (a.north east) -- node[lab, above left=0pt] {1: look up} (r.west);
  \draw[req] (a.south east) -- node[lab, below left=0pt] {2: miss} (d.west);
}

\vspace{4mm}
\panel{(C1) Results pre-computed into\\Redis before release}{%
  \node[stu] (s) {Students};
  \node[app, right=of s] (a) {App};
  \node[kv, right=of a] (r) {Redis};
  \node[db, above=5mm of r, minimum height=6.5mm, fill=orange!6] (d) {back office};
  \draw[req] (s) -- (a); \draw[req] (a) -- (r);
  \draw[pre] (d) -- node[lab, left=1pt] {before release} (r);
}\hfill
\panel{(C2) One pre-computed file per\\student, served by Nginx}{%
  \node[stu] (s) {Students};
  \node[app, right=of s] (n) {Nginx};
  \node[files, right=of n] (f) {JSON\\files};
  \node[db, above=5mm of f, minimum height=6.5mm, fill=orange!6] (d) {back office};
  \draw[req] (s) -- (n); \draw[req] (n) -- (f);
  \draw[pre] (d) -- node[lab, left=1pt] {before release} (f);
}\hfill
\panel{(D) Waiting room in front\\of A or C1}{%
  \node[stu] (s) {Students};
  \node[wr, right=8mm of s] (w) {Waiting\\room};
  \node[app, right=of w] (b) {A or C1};
  \draw[req] ([yshift=1.5mm]s.east) -- node[lab, above=1pt] {ticket} ([yshift=1.5mm]w.west);
  \draw[pre] ([yshift=-1.5mm]w.west) -- node[lab, below=1pt] {position} ([yshift=-1.5mm]s.east);
  \draw[req] (w) -- node[lab, above=1pt] {admit} (b);
}

\caption{How each design answers a request. Solid arrows are on the request
path; dashed arrows happen before release (C1, C2) or tell a queued student
their position (D).}
\label{fig:designs}
\end{figure*}

All application designs run the same Go server, which works on at most 256
requests at once and lets 4096 more wait, as real servers
do~\cite{apache2026mpm}. Real portals run full web frameworks; to stand in
for that work, building one page from the database costs \SI{5}{ms} of CPU.
This value is an assumption, which we vary in Section~\ref{sec:rq1}.

\subsection{Students and surge}

\begin{figure}[t]
\centering
% How a simulated student behaves (Section IV, Workload). Included by sections/method.tex.
\begin{tikzpicture}[
  node distance=5mm and 6mm,
  blk/.style={draw, rounded corners=2pt, align=center, font=\scriptsize, inner sep=2.5pt, fill=white},
  dec/.style={blk, fill=yellow!20},
  arr/.style={-{Stealth[length=1.6mm]}, semithick},
  lab/.style={font=\tiny, inner sep=1pt, align=center},
]
  \node[blk, fill=gray!12] (arrive) {Student arrives};
  \node[blk, fill=blue!8, below=of arrive] (ask) {Asks for result};
  \node[blk, fill=gray!6, left=9mm of ask] (early) {Before release:\\reload every\\5--15\,s};
  \node[dec, below=of ask] (what) {Answer?};
  \node[blk, fill=green!12, left=9mm of what] (ok) {Result: done\\{\tiny 10\% look again}};
  \node[blk, right=9mm of what] (bad) {Error, or nothing\\within patience\\{\tiny 5--20\,s}};
  \node[blk, fill=orange!12, below=of bad] (refresh) {Refresh after 1--5\,s\\{\tiny $p=0.8$ after error}\\{\tiny $p=0.9$ after waiting}};
  \node[blk, fill=red!10, below=of what] (quit) {Gives up\\{\tiny or after 20 tries}};
  \draw[arr] (arrive) -- (ask);
  \draw[arr] (ask.west) -- node[lab, above] {not released} (early.east);
  \draw[arr] (early.north) |- ([yshift=-0.5mm]arrive.west);
  \draw[arr] (ask) -- (what);
  \draw[arr] (what) -- (ok);
  \draw[arr] (what) -- (bad);
  \draw[arr] (bad) -- node[lab, right] {refresh} (refresh);
  \draw[arr] (bad.south west) -- node[lab, above left] {no refresh} (quit.north east);
  \draw[arr] (refresh.east) -- ++(3mm,0) |- (ask.east);
\end{tikzpicture}

\caption{How each simulated student behaves. All numbers are assumptions;
Section~\ref{sec:rq4} varies them.}
\label{fig:student}
\end{figure}

The data are \num{200000} synthetic students with six subjects each,
generated from a seed so that the right answer for any roll number can be
recomputed; no real student data was used. Students arrive as an open-loop
Poisson process~\cite{schroeder2006open}: the rate rises over the \SI{20}{s}
before release to 20\% of the peak, reaches the peak within \SI{3}{s} of
release and decays with a time constant of \SI{30}{s} over \SI{90}{s}. Real
result days stay busy far longer~\cite{ht2021ssc}; we compress the surge to
keep runs short. Early students reload every 5--\SI{15}{s} until release. Each
student waits 5--\SI{20}{s} for an answer; after an error they refresh with
probability 0.8, after waiting in vain with probability 0.9, after a pause of
1--\SI{5}{s}, for at most 20 attempts. These numbers are assumptions;
Section~\ref{sec:rq4} varies them.

\subsection{Testbed and measurement}

All runs used one desktop computer (AMD Ryzen~7 5800X, 8 cores, \SI{32}{GB},
Windows~11 with WSL\,2 and Docker~29.9). The design under test was pinned to
the two hardware threads of one core and the load generator to twelve threads
on six other cores. We first found design A's capacity: the highest constant
arrival rate at which 99\% of students got a correct result within
\SI{2}{s}, \CapA\ students/s. Loads are multiples of it; the JEE 2016
incident was about 5$\times$~\cite{pagalguy2016jee}. An oracle classifies every
response as correct, stale (the old version after a correction), wrong,
missing (``not found'' for a real student) or failed. We report the share of
students who got a correct result within 5 or \SI{10}{s} of their first
request after release, or at all, and the CPU-seconds used by all containers
per \num{1000} students served. Each configuration ran 3 times in random
order; we report medians with 95\% bootstrap intervals. Against k6's
constant-arrival-rate executor~\cite{k6docs} at 150 requests/s, our generator
measured median latencies of 5.3--\SI{5.5}{ms} against \SI{5.4}{ms}. Six
waiting-room runs in which the generator ran out of local ports were discarded
and run again with a wider port range.
